% GENERATED FILE. Edit canonical lessons, then run npm run generate:books.
% canonical-source-hash: fnv1a64:47895d3be7ed8c84
% canonical-lessons: TA-C495-cicu, TA-C495-malalai, TA-C495-taval, TA-C495-talattu, TA-C495-uttu

\chapter{Babies}
\label{ch:ta-babies}
\hlchaptermodality{495}

\begin{chapteropening}
\textbf{By the end of this chapter:} \emph{I can talk about infants, toddlers, crawling, lullabies and feeding a baby.}

Complete the last lesson of chapter 495: I can talk about infants, toddlers, crawling, lullabies and feeding a baby.
\end{chapteropening}

\section[\texorpdfstring{\ta{சிசு} (cicu)}{சிசு (cicu)}]{\ta{சிசு} (cicu) — an infant}

\label{lesson:TA-C495-cicu}



\begin{quote}
Before the new one: say the Tamil for offspring, then the Tamil for an heir.
\end{quote}

\subsection*{You'll want to know: \ta{சிசு}}
\textbf{\ta{சிசு}} — \emph{cicu} — \textquotedblleft{}an infant\textquotedblright{}.

\begin{quote}
\ta{சிசு} \ta{அமைதியாகத்} \ta{தூங்கியது}.
\end{quote}

\begin{quote}
\emph{cicu amaitiyākat tūṅkiyatu.}
\end{quote}

\begin{quote}
The infant slept peacefully.
\end{quote}

Written and medical; at home \ta{குழந்தை} or \ta{பாப்பா}.

\begin{cousinweb}
\textbf{\ta{சிசு}} (\emph{cicu}) is borrowed from Sanskrit \emph{śiśu}, \textquotedblleft{}a child\textquotedblright{} (CDIAL 12476).
\end{cousinweb}

\begin{grammarlens}[title={What you've built}]
One more word for describing what you lived and felt.
\end{grammarlens}

\subsection*{Your turn}
\begin{itemize}
\raggedright
  \item \emph{Say it:} \emph{cicu}
  \item \emph{Say it:} \emph{cicu}, once more
  \item \emph{Say it:} \emph{cicu amaitiyākat tūṅkiyatu.} — The infant slept peacefully.
  \item \emph{Recall:} say the Tamil for offspring, then the Tamil for an heir, then say \emph{cicu} again
\end{itemize}

\begin{tcolorbox}[breakable,colback=teal!4,colframe=teal!35!black,title={Before you move on}]
What does \ta{சிசு} mean? (\textquotedblleft{}An infant\textquotedblright{}.) Say it once more.
\end{tcolorbox}
\section[\texorpdfstring{\ta{மழலை} (maḻalai)}{மழலை (maḻalai)}]{\ta{மழலை} (maḻalai) — a toddler; baby talk}

\label{lesson:TA-C495-malalai}



\begin{quote}
Before the new one: say the Tamil for an heir, then the Tamil for an infant.
\end{quote}

\subsection*{You'll want to know: \ta{மழலை}}
\textbf{\ta{மழலை}} — \emph{maḻalai} — \textquotedblleft{}a toddler; baby talk\textquotedblright{}.

\begin{quote}
\ta{குழந்தையின்} \ta{மழலைப்} \ta{பேச்சு} \ta{எல்லோருக்கும்} \ta{பிடித்தது}.
\end{quote}

\begin{quote}
\emph{kuḻantaiyiṉ maḻalaip pēccu ellōrukkum piṭittatu.}
\end{quote}

\begin{quote}
Everyone loved the baby’s first words.
\end{quote}

\ta{மழலைப்} \ta{பேச்சு}, a baby’s talk; \ta{மழலையர்} \ta{பள்ளி}, a nursery school.

\begin{cousinweb}
\textbf{\ta{மழலை}} (\emph{maḻalai}) is native Dravidian (DEDR 4747), where it means childhood and tender age; the sense of a small child's talk is the everyday one now.
\end{cousinweb}

\begin{grammarlens}[title={What you've built}]
One more word for describing what you lived and felt.
\end{grammarlens}

\subsection*{Your turn}
\begin{itemize}
\raggedright
  \item \emph{Say it:} \emph{maḻalai}
  \item \emph{Say it:} \emph{maḻalai}, once more
  \item \emph{Say it:} \emph{kuḻantaiyiṉ maḻalaip pēccu ellōrukkum piṭittatu.} — Everyone loved the baby’s first words.
  \item \emph{Recall:} say the Tamil for an heir, then the Tamil for an infant, then say \emph{maḻalai} again
\end{itemize}

\begin{tcolorbox}[breakable,colback=teal!4,colframe=teal!35!black,title={Before you move on}]
What does \ta{மழலை} mean? (\textquotedblleft{}A toddler; baby talk\textquotedblright{}.) Say it once more.
\end{tcolorbox}
\section[\texorpdfstring{\ta{தவழ்} (tavaḻ)}{தவழ் (tavaḻ)}]{\ta{தவழ்} (tavaḻ) — to crawl}

\label{lesson:TA-C495-taval}



\begin{quote}
Before the new one: say the Tamil for an infant, then the Tamil for a toddler.
\end{quote}

\subsection*{You'll want to know: \ta{தவழ்}}
\textbf{\ta{தவழ்}} — \emph{tavaḻ} — \textquotedblleft{}to crawl\textquotedblright{}.

\begin{quote}
\ta{எட்டு} \ta{மாதத்தில்} \ta{குழந்தை} \ta{தவழ்ந்தது}.
\end{quote}

\begin{quote}
\emph{eṭṭu mātattil kuḻantai tavaḻntatu.}
\end{quote}

\begin{quote}
At eight months the baby crawled.
\end{quote}

Past (I …): \ta{தவழ்ந்தேன்}, \emph{tavaḻntēṉ}.

Past with -\ta{ந்த}-: \ta{தவழ்ந்தது}.

\begin{cousinweb}
\textbf{\ta{தவழ்}} (\emph{tavaḻ}) is native Dravidian (DEDR 3109), of infants, lizards and the creep of waves.
\end{cousinweb}

\begin{grammarlens}[title={What you've built}]
One more word for describing what you lived and felt.
\end{grammarlens}

\subsection*{Your turn}
\begin{itemize}
\raggedright
  \item \emph{Say it:} \emph{tavaḻ}
  \item \emph{Say it:} \emph{tavaḻ}, once more
  \item \emph{Say it:} \emph{eṭṭu mātattil kuḻantai tavaḻntatu.} — At eight months the baby crawled.
  \item \emph{Recall:} say the Tamil for an infant, then the Tamil for a toddler, then say \emph{tavaḻ} again
\end{itemize}

\begin{tcolorbox}[breakable,colback=teal!4,colframe=teal!35!black,title={Before you move on}]
What does \ta{தவழ்} mean? (\textquotedblleft{}To crawl\textquotedblright{}.) Say it once more.
\end{tcolorbox}
\section[\texorpdfstring{\ta{தாலாட்டு} (tālāṭṭu)}{தாலாட்டு (tālāṭṭu)}]{\ta{தாலாட்டு} (tālāṭṭu) — to sing a lullaby}

\label{lesson:TA-C495-talattu}



\begin{quote}
Before the new one: say the Tamil for a toddler, then the Tamil for to crawl.
\end{quote}

\subsection*{You'll want to know: \ta{தாலாட்டு}}
\textbf{\ta{தாலாட்டு}} — \emph{tālāṭṭu} — \textquotedblleft{}to sing a lullaby\textquotedblright{}.

\begin{quote}
\ta{பாட்டி} \ta{குழந்தையைத்} \ta{தாலாட்டினாள்}.
\end{quote}

\begin{quote}
\emph{pāṭṭi kuḻantaiyait tālāṭṭiṉāḷ.}
\end{quote}

\begin{quote}
Grandmother sang the baby to sleep.
\end{quote}

Past (I …): \ta{தாலாட்டினேன்}, \emph{tālāṭṭiṉēṉ}.

As a noun, a lullaby: \ta{தாலாட்டுப்} \ta{பாடு}.

\begin{cousinweb}
The origin of \textbf{\ta{தாலாட்டு}} (\emph{tālāṭṭu}) is not settled; it ends in \ta{ஆட்டு} \textquotedblleft{}to rock, to move\textquotedblright{}, but the first part is not confirmed by our sources.
\end{cousinweb}

\begin{grammarlens}[title={What you've built}]
One more word for describing what you lived and felt.
\end{grammarlens}

\subsection*{Your turn}
\begin{itemize}
\raggedright
  \item \emph{Say it:} \emph{tālāṭṭu}
  \item \emph{Say it:} \emph{tālāṭṭu}, once more
  \item \emph{Say it:} \emph{pāṭṭi kuḻantaiyait tālāṭṭiṉāḷ.} — Grandmother sang the baby to sleep.
  \item \emph{Recall:} say the Tamil for a toddler, then the Tamil for to crawl, then say \emph{tālāṭṭu} again
\end{itemize}

\begin{tcolorbox}[breakable,colback=teal!4,colframe=teal!35!black,title={Before you move on}]
What does \ta{தாலாட்டு} mean? (\textquotedblleft{}To sing a lullaby\textquotedblright{}.) Say it once more.
\end{tcolorbox}
\section[\texorpdfstring{\ta{ஊட்டு} (ūṭṭu)}{ஊட்டு (ūṭṭu)}]{\ta{ஊட்டு} (ūṭṭu) — to feed}

\label{lesson:TA-C495-uttu}



\begin{quote}
Before the new one: say the Tamil for to crawl, then the Tamil for to sing a lullaby.
\end{quote}

\subsection*{You'll want to know: \ta{ஊட்டு}}
\textbf{\ta{ஊட்டு}} — \emph{ūṭṭu} — \textquotedblleft{}to feed\textquotedblright{}.

\begin{quote}
\ta{அம்மா} \ta{குழந்தைக்குச்} \ta{சோறு} \ta{ஊட்டினாள்}.
\end{quote}

\begin{quote}
\emph{ammā kuḻantaikkuc cōṟu ūṭṭiṉāḷ.}
\end{quote}

\begin{quote}
Mother fed the child rice.
\end{quote}

Past (I …): \ta{ஊட்டினேன்}, \emph{ūṭṭiṉēṉ}.

The one fed takes -\ta{க்கு}: \ta{குழந்தைக்கு} \ta{ஊட்டு}.

\begin{cousinweb}
\textbf{\ta{ஊட்டு}} (\emph{ūṭṭu}) is native Dravidian, the causative of \ta{உண்} \emph{uṇ}, \textquotedblleft{}to eat\textquotedblright{} (DEDR 600): to make someone eat.
\end{cousinweb}

\begin{grammarlens}[title={What you've built}]
One more word for describing what you lived and felt.
\end{grammarlens}

\subsection*{Your turn}
\begin{itemize}
\raggedright
  \item \emph{Say it:} \emph{ūṭṭu}
  \item \emph{Say it:} \emph{ūṭṭu}, once more
  \item \emph{Say it:} \emph{ammā kuḻantaikkuc cōṟu ūṭṭiṉāḷ.} — Mother fed the child rice.
  \item \emph{Recall:} say the Tamil for to crawl, then the Tamil for to sing a lullaby, then say \emph{ūṭṭu} again
\end{itemize}

\begin{tcolorbox}[breakable,colback=teal!4,colframe=teal!35!black,title={Before you move on}]
What does \ta{ஊட்டு} mean? (\textquotedblleft{}To feed\textquotedblright{}.) Say it once more.
\end{tcolorbox}
